\section{The source, the date, and prior work}

William of Malmesbury reports the flight in his \textit{Gesta Regum Anglorum}
(about 1125), within living memory of the abbey's older monks and roughly a century
after the event \autocite{williamMalmesbury_gesta}.
He names Eilmer, describes the apparatus as wings on the hands and feet, gives the
distance as more than a stadium (a furlong, about \SI{200}{\meter}), reports that
wind or turbulence and his own fear brought him down, and that he broke both legs.
He adds Eilmer's retrospective explanation that he had forgotten to fit a tail.
This is the canonical account, and Lynn White's study remains the standard
historical treatment \autocite{white1961eilmer}; later popular and survey histories
follow it \autocite{hallion2003takingflight,scott1995shoulders,lacey2004greattales},
and the \textit{Oxford Dictionary of National Biography} dates Eilmer to about
985 and after 1066 \autocite{thomson2004eilmer_odnb}.

The date of the flight is an open and actively debated question.
The traditional reconstruction ties Eilmer's later comet sighting to the returns of
Halley's comet and places the flight around 1000--1010
\autocite{woosnam1986eilmer,lewispz_bayeuxcomet}.
\textcite{aitcheson2026date}, writing this year in \textit{Notes and Queries},
reopens it: the earlier comet Eilmer recalled may have been the comet of 1018 rather
than the 989 apparition of Halley, which would move his birth into the 990s to 1010s
and the flight itself as late as the 1020s, 1030s, or early 1040s, decades later than
the standard date.
The window is therefore wider and less settled than the textbooks imply, and we do
not try to close it.
Because the flight depends on weather rather than on the calendar year, we treat the
launch day as a variable sampled from the wind climatology of the site
(\cref{sec:results}).
That choice makes our analysis insensitive to the dating dispute: the relevant
prevailing-wind regime is stable across the whole eleventh-century window Aitcheson
admits, so the feasibility verdict holds whichever decade the leap belongs to.

Human-flight attempts before the modern era are mostly legend, and the
historiography is careful to separate reportable fact from myth-making
\autocite{paz2013humanflight}.
Eilmer is unusual in having a credible near-contemporary witness and a named,
physical cause of failure.
The closest antecedent is the ninth-century Andalusian polymath Abbas Ibn Firnas,
who is reported to have glided and likewise landed badly, traditionally also blamed
on the lack of a tail \autocite{anderson2024bridge}.
That two of the earliest credible accounts converge on the same aerodynamic
diagnosis is itself worth a quantitative test.
The name of the modern open-source hypersonic flow solver \textit{Eilmer} honours
this same monk \autocite{gibbons2022eilmer,jacobs2017eilmer4}, a small reminder that
the aerodynamics community already claims him.
